% Generated draft: requires mathematical and source review.
\documentclass[11pt]{article}
\usepackage[margin=25mm]{geometry}
\usepackage{fontspec}
\setmainfont{Latin Modern Roman}
\usepackage{amsmath,amssymb,mathtools,unicode-math}
\setmathfont{Latin Modern Math}
\usepackage{enumitem,needspace,xurl,hyperref}
\setlength{\parindent}{0pt}
\setlength{\parskip}{5pt}
\setlength{\emergencystretch}{3em}
\newcommand{\N}{\mathbb{N}}
\newcommand{\Z}{\mathbb{Z}}
\newcommand{\Q}{\mathbb{Q}}
\newcommand{\R}{\mathbb{R}}
\newcommand{\C}{\mathbb{C}}
\newcommand{\F}{\mathbb{F}}
\newcommand{\A}{\mathbb{A}}
\newcommand{\PP}{\mathbb{P}}
\newcommand{\E}{\mathbb{E}}
\newcommand{\eps}{\varepsilon}
\newcommand{\dd}{\,\mathrm{d}}
\begin{document}
\section*{UnsolvedMath ID: 700995}
\textbf{Source problem code:} catalogue entry
\subsection{Definitions and mathematical statement}
Write $(\Q,+)$ for the additive group of rational numbers. For a positive integer $n$, let
\[
GL_n(\Q)=\{A\in M_n(\Q):\det A\ne 0\},
\]
with matrix multiplication. An embedding is an injective group homomorphism. A group is finitely generated if some finite subset generates it, and finitely presented if it admits a presentation $\langle S\mid R\rangle$ with both $S$ and $R$ finite.

\textbf{Part (c), the remaining task in the supplied record.} Given a positive integer $n$, find an explicit and ``natural'' finitely presented group $\Gamma_n$ together with an explicit embedding
\[
\iota_n:GL_n(\Q)\hookrightarrow\Gamma_n.
\]
The source does not specify the range or quantification of $n$. The formulation above uses the conventional domain of positive matrix sizes and treats $n$ as a parameter; whether a construction is required for every positive integer $n$, or uniformly in $n$, requires confirmation. The words ``explicit'' and ``natural'' have no supplied mathematical criteria, so an arbitrary existence statement does not capture all of the stated request.

\textbf{Archived part (a).} Find an explicit and ``natural'' finitely presented group $\Gamma$ and an embedding
\[
\iota:(\Q,+)\hookrightarrow\Gamma.
\]
The source calls this a well-known problem and recalls, as background, the assertion that every recursively presented group embeds in a finitely presented group, attributed to G. Higman. This existence background does not specify the requested explicit and ``natural'' construction. The supplied archive reports that part (a) has been solved.

\textbf{Archived part (b).} Find a finitely generated group $G$ and an explicit embedding
\[
 j:(\Q,+)\hookrightarrow G.
\]
Finite presentation and ``naturalness'' are not required in the wording of this part. The source attributes the existence of such a group to Theorem IV of the cited paper of G. Higman, B. H. Neumann and H. Neumann, and reports that an explicit embedding has subsequently been found.

\textbf{Supplied covering-space phrasing.} The source also asks for a simplicial complex $X$ covering a finite complex, with fundamental group isomorphic to $(\Q,+)$ or, respectively, $GL_n(\Q)$. A precise conventional interpretation is to seek connected simplicial complexes $X,Y$, with $Y$ finite, a covering map $p:X\to Y$, and a basepoint $x\in X$, such that
\[
\pi_1(X,x)\cong(\Q,+)
\quad\text{or, respectively,}\quad
\pi_1(X,x)\cong GL_n(\Q).
\]
The covering is not required to have finitely many sheets. Connectedness and the simplicial interpretation of the finite base complex are conventions used here to make the phrasing precise; they are not explicit in the source. This formulation expresses the underlying embedding into a finitely presented fundamental group relevant to parts (a) and (c). It does not by itself specify ``explicit'' or ``natural,'' and its relationship to the weaker finite-generation requirement in part (b) is not clarified by the source.

The supplied registry labels the whole record partially solved, with parts (a) and (b) archived as solved and part (c) remaining in the main section. These are supplied source-status reports, not an independent verification of the constructions or their current literature status.
\subsection{Short English statement}
The remaining request is to construct an explicit, ``natural'' finitely presented group containing $GL_n(\Q)$ by an explicit embedding. The intended range of $n$ and the meaning of ``natural'' are unspecified. Related requests to embed the additive group $(\Q,+)$ into an explicit finitely presented group, or explicitly into a finitely generated group, are reported solved in the supplied archive. The record also gives a covering-space formulation, whose scope and additional explicitness requirements need review.
\subsection{Sources}
\begin{itemize}
\item Local UnsolvedMath registry, permanent ID $700995$, exact source problem code \texttt{catalogue entry}, title \emph{Kourovka Notebook Problem 14.10}; supplied complete record, upstream revision \texttt{d371ee917ac9647fbb06cd6fc529deec6f3dda82}.
\item \emph{The Kourovka Notebook (Unsolved Problems in Group Theory)}, problem 14.10, proposed by P. de la Harpe in the 14th issue (1999). The supplied background identifies an update in the 21st issue (2026), arXiv:1401.0300v48. The supplied collection description names E. I. Khukhro and V. D. Mazurov as editors of Issues 1--21.
\item G. Higman, \emph{Proc. Royal Soc. London Ser. A} \textbf{262} (1961), 455--475. Supplied reference for the background embedding theorem concerning recursively presented groups; no article title is supplied.
\item G. Higman, B. H. Neumann and H. Neumann, \emph{J. London Math. Soc.} \textbf{24} (1949), 247--254, Theorem IV. Supplied reference for existence in part (b); no article title is supplied.
\item J. Belk, J. Hyde and F. Matucci, \emph{Bull. Amer. Math. Soc.} \textbf{59}, no. 4 (2022), 561--567. The supplied Archive of Solved Problems entry reports that the embedding requested in part (a) is found there; no article title is supplied.
\item V. H. Mikaelian, \emph{Int. J. Math. Math. Sci.} \textbf{2005}, no. 13 (2005), 2119--2123. The supplied Archive of Solved Problems entry reports that the embedding requested in part (b) is found there; no article title is supplied.
\item The supplied status review is dated \texttt{2026-10-16}, records evidence as ``published,'' and labels the record \texttt{partially\_solved}. This assessment is supplied dated curation, not independently verified. Its note identifies parts (a) and (b) as archived solved parts and the remaining part as belonging to the main section. The Mikaelian citation is truncated in that note; the complete bibliographic details above come from the supplied background.
\end{itemize}
\end{document}
